%% ma: L12-B08-0107
%% lop: 12
%% chuong: 2
%% bai: 8
%% dang: 12.8.27
%% loai: TN4
%% muc: NB
%% dap_an: "C"
%% nguon: "CĐ2.3.pdf giáo viên gửi 27/10/2026 (M07, Lê Minh Khôi – Nguyễn Anh Khôi, Chương II Chủ đề 3 Biểu thức tọa độ của các phép toán vectơ), Câu 18"
%% kien_thuc: "Bài 8 (tọa độ trọng tâm tam giác), Bài 7 (tọa độ điểm)"
%% trang_thai: chinh-thuc
\begin{ex}
Trong không gian $Oxyz$, cho vectơ $\overrightarrow{OA}=3\vec i-2\vec j-2\vec k$ và điểm $B(0;1;-4)$. Tọa độ trọng tâm $G$ của tam giác $OAB$ là
\choice{$G(1;-1;-2)$}{$G(-1;-1;-2)$}{\True $G\left(1;-\dfrac13;-2\right)$}{$G\left(1;-\dfrac13;-\dfrac23\right)$}
\loigiai{$A(3;-2;-2)$, $G\left(\dfrac{0+3+0}3;\dfrac{0-2+1}3;\dfrac{0-2-4}3\right)=\left(1;-\dfrac13;-2\right)$.}
\end{ex}
